\documentclass[a4paper,twoside]{article}

\usepackage{morewrites}

\usepackage[svgnames,dvipsnames,x11names]{xcolor}
\usepackage{graphicx}
\usepackage[english]{babel}
\usepackage[colorlinks=true, linkcolor=NavyBlue, urlcolor=NavyBlue, citecolor=NavyBlue]{hyperref}
\usepackage[a4paper]{geometry}
\usepackage{ts-fonts}
\usepackage{csquotes}
\usepackage{amsmath}
\usepackage{float}
\usepackage{seqsplit}
\usepackage{ts-code}

\usepackage{lastpage}
\usepackage{refcount}
\makeatletter
\def\ps@plain{
\let\@oddhead\@empty
\let\@evenhead\@empty
\def\@oddfoot{
\hfil
\ifnum\getpagerefnumber{LastPage}>1\relax
\thepage
\fi
\hfil}
\let\@evenfoot\@oddfoot
}
\makeatother

\title{Metadata Conventions}
\author{}\date{}
\begin{document}
\maketitle
TeXSmith normalizes a handful of common front matter fields so that templates and fragments can lean on a single canonical name. External configuration files, CLI \tscodeinline{-\allowbreak{}-\allowbreak{}attribute} overrides, and Markdown \tscodeinline{press.*} blocks are all merged into one flat namespace before the template resolver runs, so manifests can simply reference the final attribute (e.g. \tscodeinline{emoji}, \tscodeinline{glossary\_style}, \tscodeinline{width}) without worrying about where it originally came from. The full \tscodeinline{press} tree is still kept around for backwards compatibility, but no other part of the codebase needs to dig through dotted \tscodeinline{press.*} paths anymore.
\section{Title \& Subtitle}
\tscodeinline{title} and \tscodeinline{subtitle} do exactly what the tin says: declare them in the front matter and they become the document's title and subtitle inside the template.

\tscodeinline{title} is not strictly required. If omitted, TeXSmith falls back to the first heading in the document, on the assumption that you probably meant \emph{that} to be the title. To opt out of this fallback (and render with no title at all), set \tscodeinline{title: null} explicitly:
\begin{tscode}[lang=md, engine=pygments]
\begin{Verbatim}[commandchars=\\\{\},breaklines, breakanywhere, commandchars=\\\{\}]
\PYZhy{}\PYZhy{}\PYZhy{}
\PY{g+gu}{subtitle: \PYZdq{}An In\PYZhy{}depth Exploration\PYZdq{}}
\PY{g+gu}{\PYZhy{}\PYZhy{}\PYZhy{}}
\PY{g+gh}{\PYZsh{} This section heading will not be used as the title}
\end{Verbatim}
\end{tscode}
\begin{tscode}[lang=md, engine=pygments]
\begin{Verbatim}[commandchars=\\\{\},breaklines, breakanywhere, commandchars=\\\{\}]
\PYZhy{}\PYZhy{}\PYZhy{}
\PY{g+gu}{title: Some Title}
\PY{g+gu}{\PYZhy{}\PYZhy{}\PYZhy{}}
\PY{g+gh}{\PYZsh{} This section heading will not be used as the title}
\end{Verbatim}
\end{tscode}
\section{Authors}
Author metadata is validated with Pydantic and normalized into a list of objects shaped like \tscodeinline{\{ name, affiliation \}}. The parser is forgiving on input, several syntaxes are accepted, all of which collapse to the same canonical structure:
\begin{tscode}[lang=yaml, engine=pygments]
\begin{Verbatim}[commandchars=\\\{\},breaklines, breakanywhere, commandchars=\\\{\}]
\PY{n+nn}{\PYZhy{}\PYZhy{}\PYZhy{}}
\PY{n+nt}{authors}\PY{p}{:}\PY{+w}{ }\PY{l+s}{\PYZdq{}}\PY{l+s}{Ada}\PY{n+nv}{ }\PY{l+s}{Lovelace}\PY{l+s}{\PYZdq{}}
\PY{n+nn}{\PYZhy{}\PYZhy{}\PYZhy{}}
\end{Verbatim}
\end{tscode}
\begin{tscode}[lang=yaml, engine=pygments]
\begin{Verbatim}[commandchars=\\\{\},breaklines, breakanywhere, commandchars=\\\{\}]
\PY{n+nn}{\PYZhy{}\PYZhy{}\PYZhy{}}
\PY{n+nt}{authors}\PY{p}{:}
\PY{+w}{  }\PY{p+pIndicator}{\PYZhy{}}\PY{+w}{ }\PY{l+s}{\PYZdq{}}\PY{l+s}{Ada}\PY{n+nv}{ }\PY{l+s}{Lovelace}\PY{l+s}{\PYZdq{}}
\PY{+w}{  }\PY{p+pIndicator}{\PYZhy{}}\PY{+w}{ }\PY{l+s}{\PYZdq{}}\PY{l+s}{Grace}\PY{n+nv}{ }\PY{l+s}{Hopper}\PY{l+s}{\PYZdq{}}
\PY{n+nn}{\PYZhy{}\PYZhy{}\PYZhy{}}
\end{Verbatim}
\end{tscode}
\begin{tscode}[lang=yaml, engine=pygments]
\begin{Verbatim}[commandchars=\\\{\},breaklines, breakanywhere, commandchars=\\\{\}]
\PY{n+nn}{\PYZhy{}\PYZhy{}\PYZhy{}}
\PY{n+nt}{authors}\PY{p}{:}
\PY{+w}{  }\PY{p+pIndicator}{\PYZhy{}}\PY{+w}{ }\PY{n+nt}{name}\PY{p}{:}\PY{+w}{ }\PY{l+lScalar+lScalarPlain}{Ada}\PY{l+lScalar+lScalarPlain}{ }\PY{l+lScalar+lScalarPlain}{Lovelace}
\PY{+w}{    }\PY{n+nt}{affiliation}\PY{p}{:}\PY{+w}{ }\PY{l+lScalar+lScalarPlain}{Analytical}\PY{l+lScalar+lScalarPlain}{ }\PY{l+lScalar+lScalarPlain}{Engine}
\PY{+w}{  }\PY{p+pIndicator}{\PYZhy{}}\PY{+w}{ }\PY{n+nt}{name}\PY{p}{:}\PY{+w}{ }\PY{l+lScalar+lScalarPlain}{Grace}\PY{l+lScalar+lScalarPlain}{ }\PY{l+lScalar+lScalarPlain}{Hopper}
\PY{+w}{    }\PY{n+nt}{affiliation}\PY{p}{:}\PY{+w}{ }\PY{l+lScalar+lScalarPlain}{US}\PY{l+lScalar+lScalarPlain}{ }\PY{l+lScalar+lScalarPlain}{Navy}
\PY{n+nn}{\PYZhy{}\PYZhy{}\PYZhy{}}
\end{Verbatim}
\end{tscode}
\begin{tscode}[lang=yaml, engine=pygments]
\begin{Verbatim}[commandchars=\\\{\},breaklines, breakanywhere, commandchars=\\\{\}]
\PY{n+nn}{\PYZhy{}\PYZhy{}\PYZhy{}}
\PY{n+nt}{authors}\PY{p}{:}
\PY{+w}{  }\PY{n+nt}{name}\PY{p}{:}\PY{+w}{ }\PY{l+lScalar+lScalarPlain}{Ada}\PY{l+lScalar+lScalarPlain}{ }\PY{l+lScalar+lScalarPlain}{Lovelace}
\PY{+w}{  }\PY{n+nt}{affiliation}\PY{p}{:}\PY{+w}{ }\PY{l+lScalar+lScalarPlain}{Analytical}\PY{l+lScalar+lScalarPlain}{ }\PY{l+lScalar+lScalarPlain}{Engine}
\PY{n+nn}{\PYZhy{}\PYZhy{}\PYZhy{}}
\end{Verbatim}
\end{tscode}
Each entry is trimmed, validated, and stored as \tscodeinline{\{ "name": ..., "affiliation": ... \}}. A missing \tscodeinline{name} is a hard error: TeXSmith bails out before the template render even starts, so you never end up with a half-credited document.

The legacy singular \tscodeinline{author} key is still understood by the parser for backwards compatibility, but the plural \tscodeinline{authors} is the canonical form, prefer it whenever you can.
\section{Date}
The \tscodeinline{date} field is parsed into a real \tscodeinline{datetime} object, which means templates can apply the \tscodeinline{date} filter to format it however they please. Several input shapes are accepted:
\begin{tscode}[lang=yaml, engine=pygments]
\begin{Verbatim}[commandchars=\\\{\},breaklines, breakanywhere, commandchars=\\\{\}]
\PY{n+nn}{\PYZhy{}\PYZhy{}\PYZhy{}}
\PY{n+nt}{date}\PY{p}{:}\PY{+w}{ }\PY{l+lScalar+lScalarPlain}{2024\PYZhy{}07\PYZhy{}01}\PY{+w}{ }\PY{c+c1}{\PYZsh{} ISO format}
\PY{n+nt}{date}\PY{p}{:}\PY{+w}{ }\PY{l+s}{\PYZdq{}}\PY{l+s}{Custom}\PY{n+nv}{ }\PY{l+s}{date}\PY{n+nv}{ }\PY{l+s}{string}\PY{l+s}{\PYZdq{}}
\PY{n+nt}{date}\PY{p}{:}
\PY{+w}{  }\PY{n+nt}{year}\PY{p}{:}\PY{+w}{ }\PY{l+lScalar+lScalarPlain}{2024}
\PY{+w}{  }\PY{n+nt}{month}\PY{p}{:}\PY{+w}{ }\PY{l+lScalar+lScalarPlain}{7}
\PY{+w}{  }\PY{n+nt}{day}\PY{p}{:}\PY{+w}{ }\PY{l+lScalar+lScalarPlain}{1}
\PY{n+nt}{date}\PY{p}{:}\PY{+w}{ }\PY{l+lScalar+lScalarPlain}{commit}
\end{Verbatim}
\end{tscode}
The special \tscodeinline{commit} value resolves to the date of the most recent Git commit touching the document, perfect for stamping a \enquote{last updated} date without ever editing it by hand.

Once rendered, the date is formatted according to the template's locale. \tscodeinline{2024-\allowbreak{}07-\allowbreak{}01} lands as \enquote{July 1, 2024} in an English template, or \enquote{1 de julio de 2024} in a Spanish one. \LaTeX{} (with a little help from \tscodeinline{babel}) handles the linguistic gymnastics.
\section{Version}
The \tscodeinline{version} field is a free-form label describing where the document sits in its lifecycle. It is intentionally unopinionated, anything goes: semver, calendar versioning, a Git ref, or a human-readable tag like \enquote{Confidential Draft}.
\begin{tscode}[lang=yaml, engine=pygments]
\begin{Verbatim}[commandchars=\\\{\},breaklines, breakanywhere, commandchars=\\\{\}]
\PY{n+nn}{\PYZhy{}\PYZhy{}\PYZhy{}}
\PY{n+nt}{version}\PY{p}{:}\PY{+w}{ }\PY{l+lScalar+lScalarPlain}{Confidential}\PY{l+lScalar+lScalarPlain}{ }\PY{l+lScalar+lScalarPlain}{Draft}\PY{+w}{ }\PY{c+c1}{\PYZsh{} Arbitrary string}
\PY{n+nt}{version}\PY{p}{:}\PY{+w}{ }\PY{l+lScalar+lScalarPlain}{1.0.0}\PY{+w}{ }\PY{c+c1}{\PYZsh{} Semver format}
\PY{n+nt}{version}\PY{p}{:}
\PY{+w}{  }\PY{n+nt}{major}\PY{p}{:}\PY{+w}{ }\PY{l+lScalar+lScalarPlain}{1}
\PY{+w}{  }\PY{n+nt}{minor}\PY{p}{:}\PY{+w}{ }\PY{l+lScalar+lScalarPlain}{0}
\PY{+w}{  }\PY{n+nt}{patch}\PY{p}{:}\PY{+w}{ }\PY{l+lScalar+lScalarPlain}{0}
\PY{n+nt}{version}\PY{p}{:}\PY{+w}{ }\PY{l+lScalar+lScalarPlain}{git}
\end{Verbatim}
\end{tscode}
The special \tscodeinline{git} value resolves to the latest Git tag (falling back to a short commit hash if no tag exists), giving you automatic, repository-driven versioning with zero manual upkeep.
\section{Fragments' Metadata}
Fragments, the small, composable extensions that plug into a \LaTeX{} template, can declare their own metadata schema in a \tscodeinline{fragment.toml} file. Their attributes get merged into the same flat namespace as everything else, so a fragment-defined key looks no different from a built-in one at render time. See the Fragment Guide for the full declaration syntax and resolution rules.
\end{document}
